\documentclass[11pt]{article}

\usepackage{fullpage}
\usepackage{amsmath}
\usepackage{amssymb}
\usepackage{amsfonts}
\usepackage{mathtools}
\usepackage{enumitem}
\usepackage[colorlinks,linkcolor=magenta,citecolor=blue]{hyperref}

\newcommand{\Rbar}{\bar R}
\newcommand{\E}{\mathbb E}
\newcommand{\Normal}{\mathcal N}
\newcommand{\Unif}{\operatorname{Unif}}
\newcommand{\diag}{\operatorname{diag}}

\title{Ideal-PIP Relaxed-Column Discussion}
\author{}
\date{}

\begin{document}
\maketitle

\section{Goal}

We want a principled hierarchical model for fine-mapping with out-of-sample
LD that has both of the following properties.
\begin{enumerate}[leftmargin=*]
\item When \(L=1\), the single-effect PIP is the ideal SER PIP. In particular,
because the in-sample LD matrix has diagonal entries equal to one, the PIP for
SNP \(j\) should depend on the active summary statistic \(x_j\), not on the
unknown off-diagonal entries of the in-sample LD matrix.
\item When \(L>1\), the off-diagonal LD values should still matter through the
iterative residual updates. We want relaxed LD columns to help remove fitted
effects when the reference LD \(\Rbar\) differs from the in-sample LD.
\end{enumerate}

The current SuSiE-relax idea has the second property, but its single-effect PIP
is not exactly the ideal PIP. The prime variant makes the single-effect PIP
closer to the desired form, but empirically the PIP can become too diffuse.
The target here is to separate the two roles of LD uncertainty:
\[
\text{SNP selection evidence}
\qquad\text{versus}\qquad
\text{effect-shape estimation for residual removal}.
\]

\section{Shared-\texorpdfstring{\(R\)}{R} Model}

A natural principled starting point is the original shared-\(R\) model, with
\(R\) treated as the unknown in-sample correlation matrix:
\[
R\sim p(R\mid \Rbar),
\qquad
\diag(R)=1,
\]
\[
\gamma_\ell\sim\Unif\{1,\ldots,J\},
\qquad
\beta_\ell\sim\Normal(0,\sigma_\ell^2),
\qquad
\ell=1,\ldots,L,
\]
and
\[
b=\sum_{\ell=1}^L\beta_\ell e_{\gamma_\ell},
\qquad
x\mid R,b
\sim
\Normal(Rb,R/N).
\]
This model is attractive because \(R\) appears in both the mean and the noise
covariance. The same unknown LD matrix that shapes the signal also shapes the
sampling noise.

\section{Why the SER PIP Is Ideal}

For \(L=1\), conditional on fixed \(R\), compare an active model
\(\gamma=j\) to the null model. The likelihood ratio contribution involving
\(\beta\) is
\[
\exp\left\{
N\beta R_{\cdot j}^{\top}R^{-1}x
-
\frac{N}{2}\beta^2 R_{\cdot j}^{\top}R^{-1}R_{\cdot j}
\right\}.
\]
Since
\[
R_{\cdot j}^{\top}R^{-1}x=x_j,
\qquad
R_{\cdot j}^{\top}R^{-1}R_{\cdot j}=R_{jj}=1,
\]
the likelihood ratio simplifies to
\[
\exp\left\{
N\beta x_j-\frac{N}{2}\beta^2
\right\}.
\]
Thus the active-SNP evidence depends on \(R\) only through the diagonal entry
\(R_{jj}\), which is one. Integrating over
\(\beta\sim\Normal(0,\sigma_0^2)\) gives the usual SER approximate Bayes
factor:
\[
\operatorname{ABF}_j
=
(1+N\sigma_0^2)^{-1/2}
\exp\left\{
\frac{N}{2}
\frac{N\sigma_0^2}{1+N\sigma_0^2}
x_j^2
\right\}.
\]
Therefore, after also integrating over the unknown \(R\),
\[
\Pr(\gamma=j\mid x)
\propto
\operatorname{ABF}_j.
\]
All off-diagonal LD uncertainty is common nuisance evidence for the
single-effect PIP and cancels in the normalization over \(j\).

\section{A useful decomposition}

For any correlation matrix \(R\), define
\[
S_j=R_{-j,-j}-R_{-j,j}R_{j,-j}.
\]
Under the single-effect likelihood
\[
x\mid R,\beta,\gamma=j
\sim
\Normal(\beta R_{\cdot j},R/N),
\]
we have
\[
x_j\mid R,\beta,\gamma=j
\sim
\Normal(\beta,1/N),
\]
because \(R_{jj}=1\). Also, by the Gaussian conditioning formula,
\begin{align*}
x_{-j}\mid x_j,R,\beta,\gamma=j
&\sim
\Normal\left(
\beta R_{-j,j}
+
R_{-j,j}(x_j-\beta),
\frac{1}{N}S_j
\right)\\
&=
\Normal\left(
x_jR_{-j,j},
\frac{1}{N}S_j
\right).
\end{align*}
Therefore
\begin{equation}
p(x\mid R,\beta,\gamma=j)
=
\Normal(x_j\mid \beta,1/N)
\Normal\left(x_{-j}\mid x_jR_{-j,j},\frac{1}{N}S_j\right).
\end{equation}
This is the key decomposition: conditional on \(x_j\), the distribution of
the nonactive coordinates does not involve \(\beta\).

\section{Choice of distribution of correlation matrices}

Can we have a distribution of correlation matrices that is concentrated around
\(\Rbar\), with a parameter like \(\nu\) controlling the concentration? The
ideal-PIP argument does not require a specific distribution for \(R\). It only
requires
\[
R\in\mathcal C_J
=
\{R:R=R^\top,\ R\succ0,\ \diag(R)=1\},
\]
and a prior \(p(R\mid\Rbar)\) that does not depend on \(\gamma\) or \(\beta\).
Therefore the main modeling problem is choosing a useful prior on
\(\mathcal C_J\).

The prior should satisfy the following desiderata.
\begin{enumerate}[leftmargin=*]
\item It should be centered near the reference LD matrix \(\Rbar\).
\item It should have a scalar or low-dimensional concentration parameter,
such as \(\nu\) or \(\nu_0\), controlling how much the in-sample LD can deviate
from \(\Rbar\).
\item It should allow tractable, or at least locally tractable, posterior
updates for active LD columns.
\end{enumerate}

One thing we cannot have exactly is a nondegenerate distribution on
\(\mathcal C_J\) with Gaussian column marginals
\[
R_{-j,j}\sim
\Normal(\Rbar_{-j,j},\bar S_j/\nu)
\]
for all \(j\). A Gaussian vector has support on all of
\(\mathbb R^{J-1}\), whereas a correlation column is bounded and constrained
by positive definiteness. Thus Gaussian relaxed columns should be viewed as
local approximations, not as exact prior marginals of a correlation-matrix
distribution.

Several candidate priors are worth considering.

\paragraph{Standardized inverse-Wishart.}
Draw a covariance matrix and standardize it:
\[
\tilde R\sim \mathcal{IW}(\nu_0R_0,\nu_0+J+1),
\qquad
R=\tilde D^{-1/2}\tilde R\tilde D^{-1/2},
\qquad
\tilde D=\diag(\tilde R).
\]
This is a rigorous prior on correlation matrices and is close to the
inverse-Wishart calculations we have already used. For large \(\nu_0\), it
should give approximately Gaussian local fluctuations around \(R_0\). However,
the induced distribution of \(R\) is not inverse-Wishart, the column marginals
are not exactly Gaussian, and \(E(R)\) is not necessarily exactly \(R_0\).
Thus \(R_0=\Rbar\) should initially be interpreted as a centering matrix, not
as an exact prior mean unless calibrated.

\paragraph{LKJ-type prior centered at \(\Rbar\).}
The LKJ prior is a native prior on correlation matrices, with a concentration
parameter, but the standard form is centered at the identity matrix. One could
consider a transformed or tilted LKJ-type prior centered near \(\Rbar\). This
has clean support but may not lead to simple Schur-complement column updates.

\paragraph{Partial-correlation or vine parameterization.}
Correlation matrices can be parameterized by partial correlations. We could
put shrinkage priors on those partial correlations, centered at the
corresponding partial correlations of \(\Rbar\). This gives exact support and
flexible local uncertainty. The drawback is that the induced distribution on
ordinary LD columns \(R_{-j,j}\) may be algebraically awkward.

\paragraph{Matrix-log or tangent-space prior.}
Another possibility is to map correlation matrices into an unconstrained
tangent-like space, put a Gaussian prior around a transform of \(\Rbar\), and
map back to \(\mathcal C_J\). This is appealing for local perturbation theory,
but the Jacobian and the induced column conditionals may be inconvenient for
inference.

\paragraph{Projected Gaussian perturbation.}
We could perturb \(\Rbar\) by a Gaussian symmetric matrix and project the
result back to \(\mathcal C_J\). This may be computationally convenient, but
the projection creates a density and dependence structure that is hard to
interpret as a clean hierarchical model.

At this stage, the most promising strategy is:
\begin{enumerate}[leftmargin=*]
\item Use a proper distribution on \(\mathcal C_J\), such as standardized
inverse-Wishart, to make the hierarchical model rigorous.
\item Use the exact ideal-PIP identity, which is prior-agnostic as long as
\(\diag(R)=1\).
\item Use Gaussian Schur-complement column updates as local variational
approximations for computation, rather than claiming the exact prior has
Gaussian column marginals.
\end{enumerate}

\section{standardized IW}

We consider
\[
\tilde R \sim \mathcal{IW}(\nu_0 \Rbar,\nu_0+J+1),
\qquad
R=\tilde D^{-1/2}\tilde R\tilde D^{-1/2},
\qquad
\tilde D=\diag(\tilde R).
\]
Let \(p_{\mathrm{SIW}}(R\mid\Rbar,\nu_0)\) denote the induced distribution on
correlation matrices. Still consider the SER model:
\[
\gamma\sim\Unif\{1,\ldots,J\},
\qquad
\beta\sim\Normal(0,\sigma_0^2),
\]
\begin{equation}
x\mid R,\beta,\gamma=j
\sim
\Normal(\beta R_{\cdot j},R/N).
\end{equation}

\subsection*{Exact posterior factorization}

Let
\[
p_0(x\mid R)=\Normal(x\mid0,R/N).
\]
Using the likelihood-ratio identity from the ideal-PIP derivation,
\[
p(x\mid R,\beta,\gamma=j)
=
p_0(x\mid R)
\exp\left\{
N\beta x_j-\frac{N}{2}\beta^2
\right\}.
\]
Therefore
\[
p(R,\beta\mid x,\gamma=j)
\propto
p_{\mathrm{SIW}}(R\mid\Rbar,\nu_0)p_0(x\mid R)
\,
p(\beta)
\exp\left\{
N\beta x_j-\frac{N}{2}\beta^2
\right\}.
\]
Thus \(R\) and \(\beta\) factorize in the posterior conditional on
\(\gamma=j\):
\[
p(R,\beta\mid x,\gamma=j)
=
p(R\mid x)\,p(\beta\mid x_j),
\]
where
\[
p(R\mid x)
\propto
p_{\mathrm{SIW}}(R\mid\Rbar,\nu_0)
\Normal(x\mid0,R/N)
\]
is common for all \(j\), and
\[
\beta\mid x_j,\gamma=j
\sim
\Normal(m_j,v),
\qquad
v=(N+\sigma_0^{-2})^{-1},
\qquad
m_j=vNx_j
=
\frac{N\sigma_0^2}{1+N\sigma_0^2}x_j.
\]

\subsection*{PIP}

Integrating out \(\beta\) gives the usual single-effect approximate Bayes
factor
\[
\operatorname{ABF}_j
=
(1+N\sigma_0^2)^{-1/2}
\exp\left\{
\frac{N}{2}
\frac{N\sigma_0^2}{1+N\sigma_0^2}
x_j^2
\right\}.
\]
The \(R\)-dependent factor
\[
\int
p_{\mathrm{SIW}}(R\mid\Rbar,\nu_0)
\Normal(x\mid0,R/N)\,dR
\]
is common across \(j\). Therefore
\[
\Pr(\gamma=j\mid x)
=
\frac{\operatorname{ABF}_j}
{\sum_{k=1}^J\operatorname{ABF}_k}.
\]
Thus the standardized-IW prior preserves the ideal SER PIP exactly.

\subsection*{Posterior of \(R\)}

The posterior for \(R\) is rigorous but not conjugate in the standardized-IW
parameterization. In terms of \(\tilde R\),
\[
p(\tilde R\mid x)
\propto
\mathcal{IW}(\tilde R;\nu_0\Rbar,\nu_0+J+1)
|R(\tilde R)|^{-1/2}
\exp\left\{
-\frac{N}{2}x^\top R(\tilde R)^{-1}x
\right\},
\]
where
\[
R(\tilde R)=\tilde D^{-1/2}\tilde R\tilde D^{-1/2}.
\]
Equivalently,
\[
|R(\tilde R)|
=
\frac{|\tilde R|}{\prod_{k=1}^J\tilde R_{kk}},
\qquad
R(\tilde R)^{-1}
=
\tilde D^{1/2}\tilde R^{-1}\tilde D^{1/2}.
\]
The dependence of \(\tilde D\) on \(\tilde R\) breaks inverse-Wishart
conjugacy. Exact inference for \(R\) would therefore require numerical
integration, MCMC, Laplace approximation, or another approximation.

\subsection*{Column inference}

For relaxed-column use, define
\[
z_j=R_{-j,j},
\qquad
S_j=R_{-j,-j}-z_jz_j^\top.
\]
Under the common posterior \(p(R\mid x)\), the marginal posterior for the
\(j\)-th column can be written using the decomposition above:
\[
p(z_j,S_j\mid x)
\propto
p_{\mathrm{SIW},j}(z_j,S_j\mid\Rbar,\nu_0)
\Normal\left(
x_{-j}\mid x_jz_j,\frac{1}{N}S_j
\right),
\]
where \(p_{\mathrm{SIW},j}\) is the marginal prior induced by standardized IW
on the block variables \((z_j,S_j)\). The factor
\(\Normal(x_j\mid0,1/N)\) is constant with respect to \((z_j,S_j)\) and has
been dropped.

This expression is exact, but \(p_{\mathrm{SIW},j}(z_j,S_j\mid\Rbar,\nu_0)\)
does not have the same simple normal--inverse-Wishart form as the
unstandardized inverse-Wishart prior. A practical local approximation is to use
the familiar Schur-complement working prior
\[
z_j\mid S_j
\approx
\Normal(\mu_j,S_j/\nu_{\mathrm{eff}}),
\qquad
\mu_j=\Rbar_{-j,j},
\]
with \(\nu_{\mathrm{eff}}\) calibrated to the standardized-IW concentration.
Then
\[
z_j\mid S_j,x
\approx
\Normal\left(
\frac{\nu_{\mathrm{eff}}\mu_j+Nx_jx_{-j}}
{\nu_{\mathrm{eff}}+Nx_j^2},
\frac{S_j}{\nu_{\mathrm{eff}}+Nx_j^2}
\right).
\]
Equivalently, the approximate posterior mean column is
\[
\E(z_j\mid x)
\approx
\mu_j+
\frac{Nx_j}{\nu_{\mathrm{eff}}+Nx_j^2}
(x_{-j}-x_j\mu_j).
\]
The exact model supplies the ideal PIP and a proper correlation-matrix prior;
the Gaussian column update is then a local variational approximation used for
computing relaxed effect shapes.



\end{document}
